\chapter{Lesson 7: A Birthday Party}

\section{Vocabulary}

\begin{center}
\begin{tabular}{llll}
\toprule
Characters & Pinyin & Part of speech & Meaning \\
\midrule
\zh{过生日} & guò shēngrì & verb-object & to celebrate a birthday \\
\zh{聚会} & jùhuì & noun & party, get-together \\
\zh{舞会} & wǔhuì & noun & dance party \\
\zh{礼物} & lǐwù & noun & gift \\
\zh{送} & sòng & verb & to give as a gift; to see off \\
\zh{带} & dài & verb & to bring, take \\
\zh{蛋糕} & dàngāo & noun & cake \\
\zh{水果} & shuǐguǒ & noun & fruit \\
\zh{苹果} & píngguǒ & noun & apple \\
\zh{梨} & lí & noun & pear \\
\zh{西瓜} & xīguā & noun & watermelon \\
\zh{饮料} & yǐnliào & noun & beverage \\
\zh{祝} & zhù & verb & to wish \\
\zh{快乐} & kuàilè & adjective & happy \\
\zh{长寿面} & chángshòumiàn & noun & longevity noodles \\
\zh{属} & shǔ & verb & to be born in the year of \\
\zh{帮忙} & bāng máng & verb-object & to help \\
\zh{准备} & zhǔnbèi & verb & to prepare \\
\zh{还} & hái & adverb & also, in addition; still \\
\zh{着} & zhe & particle & (ongoing state) \\
\zh{进来} / \zh{出去} & jìnlai / chūqu & verb + complement & come in / go out \\
\zh{上来} / \zh{下去} & shànglai / xiàqu & verb + complement & come up / go down \\
\zh{回来} / \zh{过来} & huílai / guòlai & verb + complement & come back / come over \\
\zh{楼上} / \zh{楼下} & lóushàng / lóuxià & noun & upstairs / downstairs \\
\bottomrule
\end{tabular}
\end{center}

\section{Birthday expressions}

\begin{itemize}
  \item \zh{祝你生日快乐！} Happy birthday!
  \item \zh{今天我过生日，请你来参加我的生日聚会。} It's my birthday today; please come to
        my party. (\zh{参加} cānjiā = to attend)
  \item \zh{这是送给你的礼物。} This is a gift for you.
  \item \zh{你属什么？}---\zh{我属龙。} What's your zodiac sign?---Dragon.
\end{itemize}

\section{Grammar}

\subsection{Simple directional complements: \zh{来} and \zh{去}}

After a verb of motion, \zh{来} indicates movement \emph{toward} the speaker and
\zh{去} movement \emph{away} from the speaker.

\begin{center}
\begin{tikzpicture}
  \node[draw, circle, fill=yellow!40, minimum size=1cm, font=\small] (sp) at (0,0) {\zh{我}};
  \node[font=\scriptsize] at (0,-0.8) {speaker};
  \node[draw, rounded corners, fill=blue!8, font=\small] (a) at (-4.5,0) {\zh{他}};
  \node[draw, rounded corners, fill=blue!8, font=\small] (b) at (4.5,0) {\zh{他}};
  \draw[->, very thick, green!50!black] (a) -- node[above, font=\small] {V + \zh{来}} node[below, font=\scriptsize] {toward me} (sp);
  \draw[->, very thick, red!70!black] (sp) -- node[above, font=\small] {V + \zh{去}} node[below, font=\scriptsize] {away from me} (b);
  % up/down/in/out
  \draw[thick] (-1.6,1.4) rectangle (1.6,3.6);
  \node[font=\scriptsize, anchor=south] at (0,3.65) {a room / upstairs};
  \draw[->, thick, blue!70] (-0.4,1.7) -- node[midway, left, font=\small] {\zh{上}} (-0.4,3.3);
  \draw[->, thick, blue!70] (0.4,3.3) -- node[midway, right, font=\small] {\zh{下}} (0.4,1.7);
  \draw[->, thick, purple] (-3.4,2.9) -- node[pos=0.25, above, font=\small] {\zh{进}} (-1.0,2.9);
  \draw[->, thick, purple] (1.0,2.1) -- node[pos=0.75, above, font=\small] {\zh{出}} (3.4,2.1);
\end{tikzpicture}
\end{center}

\begin{itemize}
  \item \zh{你快进来吧！} (said from inside) Come in!
  \item \zh{他出去了。} He's gone out.
  \item \zh{妈妈，你什么时候回来？} Mom, when are you coming back?
  \item \zh{请上来。} (said upstairs) Please come up.
\end{itemize}

\subsection{Place objects: before \zh{来/去}}

\begin{important}[Position of objects]
A \emph{place} object goes between the verb and \zh{来/去}:
\zh{他回家去了}, \zh{请进教室来}, \zh{我们上楼去吧}. (Never *\zh{他回去家了}.)
A \emph{thing} object can go either before \zh{来/去} or after the whole verb phrase:
\zh{带一个蛋糕来} / \zh{带来一个蛋糕}.
\end{important}

\subsection{Compound directional complements}

Combine one of \zh{上, 下, 进, 出, 回, 过, 起} with \zh{来/去}:

\begin{center}
\begin{tabular}{llllllll}
\toprule
 & \zh{上} & \zh{下} & \zh{进} & \zh{出} & \zh{回} & \zh{过} & \zh{起} \\
\midrule
\zh{来} & \zh{上来} & \zh{下来} & \zh{进来} & \zh{出来} & \zh{回来} & \zh{过来} & \zh{起来} \\
\zh{去} & \zh{上去} & \zh{下去} & \zh{进去} & \zh{出去} & \zh{回去} & \zh{过去} & --- \\
\bottomrule
\end{tabular}
\end{center}

These follow another verb of motion: \zh{走进来} ``walk in'', \zh{跑出去} ``run out'',
\zh{站起来} ``stand up'', \zh{拿出来} ``take out''.
\begin{itemize}
  \item \zh{他从楼上走下来了。} He walked downstairs (toward me).
  \item \zh{她从包里拿出一个礼物来。} She took a gift out of her bag.
  \item \zh{大家站起来，一起唱生日歌吧！} Everyone stand up and sing ``Happy Birthday''!
\end{itemize}

\subsection{The durative particle \zh{着}}

\zh{着} after a verb describes an ongoing \emph{state} that results from an action, or an
action that accompanies another.
\begin{itemize}
  \item \zh{门开着。} The door is open (has been opened and stays open).
  \item \zh{桌子上放着一个大蛋糕。} There's a big cake on the table. (place + V\zh{着} +
        thing: existence)
  \item \zh{她穿着一条红裙子。} She's wearing a red dress.
  \item \zh{他笑着说：``生日快乐！''} Smiling, he said ``Happy birthday!'' (V1 \zh{着} V2:
        manner)
  \item \zh{别站着，坐吧！} Don't stand; sit!
\end{itemize}
Negation: \zh{门没开着。} (rare; usually just \zh{门没开}).

\subsection{\zh{还} ``also, in addition''}

\zh{还} adds another item, usually something extra: \zh{我带了蛋糕，还带了水果。}
``I brought a cake, and fruit too.'' In questions: \zh{你还要什么？} ``What else do you
want?'' (Compare \zh{也}, which adds another \emph{subject} doing the same thing.)

\section{The Chinese zodiac}

Each lunar year is assigned one of twelve animals (\zh{生肖} shēngxiào) in a fixed
cycle. The wheel shows the 2008--2019 cycle; add or subtract 12 for other years (2020 was a Rat year again).

\begin{center}
\begin{tikzpicture}
  \draw[thick] (0,0) circle (3.7);
  \draw[thick] (0,0) circle (1.4);
  \foreach \i/\a/\p/\y in {0/鼠/shǔ/2008, 1/牛/niú/2009, 2/虎/hǔ/2010, 3/兔/tù/2011, 4/龙/lóng/2012, 5/蛇/shé/2013, 6/马/mǎ/2014, 7/羊/yáng/2015, 8/猴/hóu/2016, 9/鸡/jī/2017, 10/狗/gǒu/2018, 11/猪/zhū/2019} {
    \draw[gray!60] ({90-\i*30-15}:1.4) -- ({90-\i*30-15}:3.7);
    \node[font=\large] at ({90-\i*30}:3.15) {\zh{\a}};
    \node[font=\scriptsize] at ({90-\i*30}:2.55) {\emph{\p}};
    \node[font=\scriptsize, gray!50!black] at ({90-\i*30}:1.95) {\y};
  }
  \node[align=center, font=\small] at (0,0) {\zh{十二生肖}\\ cycle of 12};
\end{tikzpicture}
\end{center}

\begin{example}[Working out a sign]
A student born in 2000 (before or after Lunar New Year matters!) is a Dragon:
2012 $-$ 12 = 2000. Someone born in 1999 (after New Year)
is a Rabbit.
\end{example}

\section{Culture: birthdays and gifts}

\begin{aside}
\begin{itemize}
  \item Traditional birthday food is \zh{长寿面} ``longevity noodles'': long uncut
        noodles for a long life. Cakes are a modern import.
  \item Gifts are given and received with both hands and often not opened in front of
        the giver.
  \item Taboo gifts rely on puns: a clock (\zh{送钟} sòng zhōng) sounds like
        \zh{送终} ``attend someone's deathbed''; pears (\zh{梨} lí) sound like
        \zh{离} ``separate''; four of anything (\zh{四} sì $\approx$ \zh{死} sǐ ``die'').
  \item In one's own zodiac year (\zh{本命年} běnmìngnián) people wear red for luck.
\end{itemize}
\end{aside}

\section{Model narrative}

\zh{昨天是小林的生日。晚上七点，我们都到他家去给他过生日。我带了一个蛋糕，小王还带来了很多水果和饮料。我们一进门，就看见桌子上放着一大碗长寿面。小林穿着一件新衬衫，笑着说：``快进来，快进来！''我们唱了生日歌，然后小林的妹妹从楼上拿下来一个大礼物。大家一直玩到十一点才回去。}

\section*{Review}
\begin{itemize}
  \item V + \zh{来} (toward speaker) / \zh{去} (away).
  \item Compound complements: \zh{上/下/进/出/回/过/起} + \zh{来/去}.
  \item Place objects go before \zh{来/去}: \zh{回家去}, \zh{走进教室来}.
  \item V\zh{着}: ongoing state; place + V\zh{着} + thing; V1\zh{着} V2.
  \item \zh{还} = in addition.
  \item Zodiac: \zh{鼠牛虎兔龙蛇马羊猴鸡狗猪}; gift taboos.
\end{itemize}
